% !TEX program = xelatex
% ============================================================
%  speech-id — Technical Report (English) — XeLaTeX
% ============================================================
\documentclass[12pt, a4paper]{article}

\usepackage{fontspec}
\usepackage{polyglossia}
\setdefaultlanguage{english}
\setmainfont{Linux Libertine O}
\setmonofont{DejaVu Sans Mono}[Scale=0.85]

\usepackage[a4paper, top=25mm, bottom=28mm, left=28mm, right=22mm]{geometry}
\usepackage{setspace}\onehalfspacing
\usepackage{parskip}\setlength{\parskip}{5pt}

\usepackage{amsmath, amssymb}
\usepackage{booktabs, tabularx, multirow, array}
\newcolumntype{R}[1]{>{\raggedleft\arraybackslash}p{#1}}
\usepackage{graphicx, float, caption}
\graphicspath{{../}{./}}
\captionsetup{font=small, labelfont=bf, skip=5pt}

\usepackage{listings, xcolor}
\definecolor{codebg}{HTML}{F5F5F5}
\definecolor{kw}{HTML}{0033B3}
\definecolor{str}{HTML}{067D17}
\definecolor{cmt}{HTML}{8C8C8C}
\lstset{basicstyle=\ttfamily\small, backgroundcolor=\color{codebg},
  frame=single, breaklines=true,
  keywordstyle=\color{kw}\bfseries, stringstyle=\color{str},
  commentstyle=\color{cmt}\itshape,
  numbers=left, numberstyle=\tiny, numbersep=7pt,
  showstringspaces=false, tabsize=4, captionpos=b,
  language=Python, morekeywords={self,True,False,None}}

\usepackage{fancyhdr}
\pagestyle{fancy}\fancyhf{}
\fancyhead[R]{\small\thepage}
\fancyhead[L]{\small\itshape Speaker Identification, speech-id}
\renewcommand{\headrulewidth}{0.4pt}

\newcommand{\code}[1]{\texttt{#1}}

\usepackage[colorlinks, linkcolor=blue!60!black, citecolor=green!40!black,
  urlcolor=blue!60!black, unicode, pdfencoding=auto]{hyperref}
\usepackage[numbers]{natbib}
\bibliographystyle{unsrtnat}
\usepackage{enumitem}
\setlist{itemsep=2pt, topsep=3pt}

% ============================================================
\begin{document}

% ── Title page ──────────────────────────────────────────────
\begin{titlepage}
\centering\vspace*{18mm}
\rule{\linewidth}{1.2pt}\\[5mm]
{\LARGE\bfseries Speaker Identification via Partition-Based Search:\\[2mm]
  A Comparative Study on VoxCeleb1}\\[5mm]
\rule{\linewidth}{1.2pt}\\[8mm]
{\large George Nazos MTN2519 Konstantinos Kaimakis MTN2508}\\[2mm]
NCSR Demokritos && UNIPI\\[2mm] June 2026}\\[14mm]

\begin{abstract}
\noindent
For our assignment in Machine Learning in Multimodal data we created \textbf{speech-id}, a speaker identification system
evaluated on VoxCeleb1 audio files that compares three vector search architectures
across three speaker embedding models.
Starting from a PostgreSQL\,+\,pgvector database of 144,840 enrollment
utterances from 1,251 speakers, we evaluate:
(i)~exhaustive brute-force scan,
(ii)~metadata hierarchy routing (gender, nationality), and
(iii)~$K$-means cluster routing (256 clusters),
each combined with three aggregation modes (speaker centroid, closest
utterance, and vote).
Experiments are conducted for ECAPA-TDNN (192\,dim),
WavLM X-Vector (512\,dim), and WavLM mean-pooled (768\,dim) embeddings
across 8,676 test utterances.
ECAPA-TDNN achieves near-perfect identification at 99.75\% Top-1 under
brute-force search, $K$-means using speaker's centroid reaches 93.56\% at only 0.015\,ms/query, a 2,520$\times$ speed-up over brute-force closest.
WavLM X-Vector yields 97.42\% Top-1 but suffers from high EER (17.9\%),
while naive WavLM mean-pooling collapses to 26--30\% Top-1, confirming
that task-specific fine-tuning dominates the choice of search strategy.
\end{abstract}
\vfill
{\footnotesize Source code:
\url{https://github.com/g-nazos/speech-id}}
\end{titlepage}

\clearpage\tableofcontents\clearpage

% ============================================================
\section{Introduction}
\label{sec:intro}

Speaker identification, determining which enrolled speaker produced
a given utterance, is a core component of multi-speaker meeting
transcription, voice-based access control, and personalised voice
assistants.
In a closed system, identifying a speaker reduces to a nearest-neighbour search against the enrolled voice data.
At database sizes of tens of thousands of utterances, the naive exhaustive search becomes a bottleneck.

This report documents an evaluation of how three search
architectures interact with three embedding models of varying
dimensionality and training objective.
All experiments share a common PostgreSQL\,+\,pgvector backend
and are run on 8,676 unseen test samples from VoxCeleb1.

Our findings are:
\begin{enumerate}
  \item A reproducible evaluation framework built on PostgreSQL\,+\,pgvector
        with separate schemas per embedding model.
  \item A three-way comparison of search strategies: brute-force,
        metadata-gated, and unsupervised cluster routing.
  \item An analysis of three aggregation modes (centroid, closest,
        vote) across all architecture--model combinations.
\end{enumerate}

% ============================================================
\section{Related Work}
\label{sec:related}

\paragraph{Speaker embeddings.}
$x$-vectors~\cite{snyder2018xvector} are fixed-length speaker embeddings
extracted by a Time Delay Neural Network (TDNN), a feed-forward
architecture that models temporal context by applying convolutions
over dilated time steps.
In the x-vector method, a statistics pooling layer aggregates the frame-level TDNN outputs
into a single utterance-level vector (mean and standard deviation),
which is then projected to a compact speaker representation.
ECAPA-TDNN extends this by adding
multi-scale Res2Net aggregation and channel-dependent frame attention~\cite{desplanques2020ecapa}, achieving state-of-the-art performance on VoxCeleb.
WavLM~\cite{chen2022wavlm} applies masked speech prediction as a
self-supervised pre-training objective over 94,000 hours of audio,
producing rich 768-dimensional representations that transfer well
to downstream tasks when fine-tuned.


% ============================================================
\section{Dataset and Split}
\label{sec:dataset}

\textbf{VoxCeleb1}~\cite{nagrani2017voxceleb} contains audio from
1,251 celebrities (690 male, 561 female) extracted from YouTube
interviews, sampled at 16\,kHz, stored as \code{.wav} files.

\paragraph{Train/test split.}
We adopt a \emph{leave-one-video-out} strategy:
\code{make\_test\_split.py} moves the lexicographically last video
directory of each speaker to the test set, minimizing the chance that test utterances come from a recording session unseen during enrollment.

\begin{table}[H]
\centering
\caption{Dataset statistics after leave-one-video-out split.}
\label{tab:dataset}
\begin{tabular}{lrr}
\toprule
\textbf{Partition} & \textbf{Utterances} & \textbf{Speakers} \\
\midrule
Enrollment (train) & 144,840 & 1,251 \\
Test               &   8,676 & 1,251 \\
\midrule
Total              & 153,516 & 1,251 \\
\bottomrule
\end{tabular}
\end{table}


% ============================================================
\section{Embedding Models}
\label{sec:embedding}

Three embedding extractors are evaluated, each stored in an isolated
PostgreSQL schema.

\subsection{ECAPA-TDNN (192-dim)}

\href{https://huggingface.co/speechbrain/spkrec-ecapa-voxceleb}%
{\code{speechbrain/spkrec-ecapa-voxceleb}} this system is composed of an ECAPA-TDNN model. It is a combination of convolutional and residual blocks ~\cite{desplanques2020ecapa,ravanelli2021speechbrain}. The model encodes a 16\,kHz waveform into a \textbf{192-dimensional}
L2-normalised embedding.
The model was already fine-tuned on VoxCeleb2 for speaker verification using
additive angular margin loss.

\subsection{WavLM X-Vector (512-dim)}

\href{https://huggingface.co/microsoft/wavlm-base-plus-sv#wavlm-base-plus-for-speaker-verification}%
{\code{wavlm-base-plus-for-speaker-verification}} is a WavLM-Base+ transformer backbone is combined with an $x$-vector head
(statistics pooling and linear projection) fine-tuned on VoxCeleb1 for speaker
verification, producing \textbf{512-dimensional} L2-normalised embeddings.
Because this model has been exposed to speaker verification objectives
during fine-tuning, it is directly comparable to ECAPA-TDNN in terms
of training signal.

\subsection{WavLM Mean-Pooled (768-dim)}

\href{https://huggingface.co/microsoft/wavlm-base-plus}%
{\code{microsoft/wavlm-base-plus}}~\cite{chen2022wavlm} is used
\emph{without any task-specific fine-tuning}: hidden states from all
12 transformer layers are mean-pooled into a \textbf{768-dimensional}
vector, clipped to $\leq$20\,s due to VRAM constraints.
This configuration tests whether generic representations
carry sufficient speaker-discriminative information for closed-set
identification without additional adaptation.

% ============================================================
\section{Database Architecture}
\label{sec:database}

\subsection{Schema Design}

PostgreSQL 16 with pgvector 0.4.2 is deployed via Docker.
Each embedding model occupies a separate schema
(\code{public} for ECAPA, \code{wavlm\_xvector}, \code{wavlm})
with identical table structure but with specific vector dimension for each model.

\begin{table}[H]
\centering
\caption{Database tables per schema. \{dim\} is model-specific (192/512/768).}
\label{tab:schema}
\begin{tabularx}{\linewidth}{llR{16mm}X}
\toprule
\textbf{Table} & \textbf{Key} & \textbf{Rows} & \textbf{Purpose} \\
\midrule
\code{speakers}            & \code{speaker\_id}   & 1,251   & Gender and nationality metadata \\
\code{audio\_embeddings}   & \code{id} (identity) & 144,840 & Per-utterance \{dim\}-dim embedding \\
\code{metadata\_centroids} & (category, value)    & 58      & Mean embeddings by gender/nationality \\
\code{cluster\_centroids}  & \code{cluster\_id}   & 256     & $K$-means cluster centres \\
\bottomrule
\end{tabularx}
\end{table}


\begin{figure}[H]
\centering
\includegraphics[width=1.00\linewidth]{db_schema.png}
\caption{Entity-relationship diagram of the database schema.
  Each embedding model uses an identical structure in its own PostgreSQL
  schema, \texttt{\{dim\}} is substituted with 192, 512, or 768 at
  schema creation time.}
\label{fig:db_schema}
\end{figure}
The \code{speakers} table is the core of the schema: it stores one
row per enrolled speaker with their gender and nationality, and every
audio embedding row carries a foreign key back to it.
The \code{audio\_embeddings} table is the largest, holding one
dimensional vector per utterance together with a
\code{cluster\_id} that records which $K$-means partition the
embedding was assigned to during the population phase.
The \code{metadata\_centroids} table stores the mean embedding
for each (category, value) pair , two gender groups and 56
gender$\times$nationality groups , so that the metadata routing step
can identify the closest partition.
Finally, \code{cluster\_centroids} stores the 256 $K$-means centres
used by the cluster-routing architecture, at query time the nearest cluster centroid is found and its \code{cluster\_id} is passed as a
\code{WHERE} filter to restrict the database scan to roughly
$144{,}840/256 \approx 566$ candidate utterances.

\subsection{Population Pipeline}

\code{populate\_db.py} runs in three phases:
\begin{enumerate}
  \item \textbf{Ingestion}: loads the \code{.pt} embeddings file and
        VoxCeleb metadata CSV, inserts speakers, streams embeddings via
        PostgreSQL binary COPY, computes 58 metadata centroids
        (2 gender + 56 gender$\times$nationality groups).
  \item \textbf{Clustering}: fetches all embeddings, L2-normalises them,
        and runs \code{MiniBatchKMeans}($K{=}256$).
  \item \textbf{Assignment}: inserts cluster centroids, updates every
        row in \code{audio\_embeddings} with its \code{cluster\_id}.
\end{enumerate}

% ============================================================
\section{Search Architectures}
\label{sec:search}

Let $\mathbf{q}\in\mathbb{R}^{d}$ be the L2-normalised query embedding.
All searches use cosine similarity
$\mathrm{sim}(\mathbf{a},\mathbf{b})=\mathbf{a}^{\top}\mathbf{b}$,
implemented via the \code{<=>} operator.

Before every evaluation run, indexed search is disabled so that
PostgreSQL performs exact sequential scans regardless of architecture:

\begin{lstlisting}[language=SQL, numbers=none,
    caption={Index scan disabled before every evaluation run.}]
SET enable_indexscan = off;
SET enable_indexonlyscan = off;
SET enable_bitmapscan = off;
\end{lstlisting}
\newpage
\subsection{Brute-Force}
\begin{lstlisting}[language=SQL, numbers=none,
    caption={Brute-force utterance-level query.}]
SELECT ae.id, ae.speaker_id,
       ae.embedding <=> %s::vector AS distance
FROM   audio_embeddings ae
JOIN   speakers s ON ae.speaker_id = s.speaker_id
ORDER  BY ae.embedding <=> %s::vector ASC
LIMIT  %s;
\end{lstlisting}

\subsection{Metadata Hierarchy}

Two in-memory comparisons route the query before issuing a filtered
SQL query:
\begin{enumerate}
  \item $g^* = \arg\max_{g}\langle\mathbf{q},\boldsymbol{\mu}_g\rangle$
        over the 2 gender centroids.
  \item Top-$P$ nationalities from gender$\times$nationality
        centroids matching $g^*$.
  \item Filtered DB query:
\end{enumerate}

\begin{lstlisting}[language=SQL, numbers=none,
    caption={Metadata-hierarchy filtered query.}]
SELECT ae.id, ae.speaker_id,
       ae.embedding <=> %s::vector AS distance
FROM   audio_embeddings ae
JOIN   speakers s ON ae.speaker_id = s.speaker_id
WHERE  s.gender = %s
  AND  s.nationality = ANY(%s)
ORDER  BY ae.embedding <=> %s::vector ASC
LIMIT  %s;
\end{lstlisting}
\newpage
\subsection{K-Means Cluster Routing}

Top-$C$ nearest cluster centroids ($C{=}1$) are identified
then only utterances in those clusters are searched:

\begin{lstlisting}[language=SQL, numbers=none,
    caption={$K$-means cluster-routed query.}]
SELECT ae.id, ae.speaker_id,
       ae.embedding <=> %s::vector AS distance
FROM   audio_embeddings ae
JOIN   speakers s ON ae.speaker_id = s.speaker_id
WHERE  ae.cluster_id = ANY(%s)
ORDER  BY ae.embedding <=> %s::vector ASC
LIMIT  %s;
\end{lstlisting}

With $K{=}256$ and 144,840 utterances, each cluster holds on average
${\approx}600$ embeddings.

\subsection{Aggregation Modes}
\begin{itemize}
    \item \textbf{Speaker centroid}: Per-speaker mean embeddings are
        loaded from the database once at startup and held in memory.
        Per-query ranking is a pure in-memory cosine dot-product:

  \item \textbf{Closest utterance}: The speaker of the single nearest
        DB result is returned.
  \item \textbf{Vote}: The top-25 DB results are numbered by
        \code{speaker\_id}, the speaker with the most votes wins, ties broken by best distance.
\end{itemize}

% ============================================================
\section{Evaluation Metrics}
\label{sec:metrics}

\begin{itemize}
  \item \textbf{Top-1\%}: fraction of test queries for which the correct speaker is ranked first, expressed as a percentage.
  \item \textbf{Hit@10\%}: fraction of queries for which the correct speaker appears anywhere in the top-10 results.
  \item \textbf{Macro-F1\%}: harmonic mean of per-speaker precision and recall, averaged equally over all 1,251 speakers.
  \item \textbf{EER\%}: decision threshold $\tau^*$ at which the false-accept rate equals the false-reject rate; lower is better.
\end{itemize}


All classification metrics are macro-averaged over 1,251 speakers.
EER is computed by sweeping the decision threshold $\tau$ on the
top-1 similarity score.
Latency is measured with \code{time.perf\_counter\_ns} and reported
as mean milliseconds per query over 8,676 test utterances.

% ============================================================
\section{Experimental Setup}
\label{sec:setup}

\begin{table}[H]
\centering
\caption{Hyperparameters shared across all model experiments.}
\label{tab:hyperparams}
\begin{tabular}{ll}
\toprule
\textbf{Parameter} & \textbf{Value} \\
\midrule
Distance metric               & Cosine (\code{<=>}) \\
Clusters $K$                  & 256 \\
Cluster probe depth $C$       & 1 \\
Metadata probe depth $P$      & 1 \\
Top-$K$ for Hit@$K$           & 10 \\
Candidate limit (vote)        & 25 \\
Test queries                  & 8,676 \\
Utterance mode                & all \\
\bottomrule
\end{tabular}
\end{table}

% ============================================================
\section{Results}
\label{sec:results}

\subsection{ECAPA-TDNN (192-dim)}
\label{sec:results_ecapa}

\begin{table}[H]
\centering
\caption{ECAPA-TDNN results (cosine, $C{=}1$, schema: \code{public}).
  \textbf{Bold}: best per column overall,
  \underline{underline}: best non-exact (non-brute-force) configuration.}
\label{tab:results_ecapa}
\setlength{\tabcolsep}{4pt}
\begin{tabular}{llrrrrr}
\toprule
\textbf{Architecture} & \textbf{Aggregation} &
\textbf{Top-1\%} & \textbf{Hit@10\%} &
\textbf{F1\%} & \textbf{EER\%} & \textbf{ms} \\
\midrule
\multirow{3}{*}{Brute-force}
  & Speaker centroid & 99.64 & \textbf{99.97} & 98.99 & 0.61 & 0.030 \\
  & Closest          & \textbf{99.75} & 99.94 & \textbf{99.06} & 0.79 & 37.81 \\
  & Vote             & 99.68 & 99.94 & 98.97 & 3.57 & 37.86 \\
\midrule
\multirow{3}{*}{Metadata hier.}
  & Speaker centroid & 33.24 & 33.29 & 28.94 & \textbf{0.17} & 0.032 \\
  & Closest          & 33.29 & 33.30 & 29.21 & 0.42 & 30.47 \\
  & Vote             & 33.25 & 33.30 & 29.12 & 0.45 & 30.21 \\
\midrule
\multirow{3}{*}{$K$-Means}
  & Speaker centroid & \underline{93.56} & \underline{93.75} & \underline{88.81} & \underline{0.90} & \textbf{0.015} \\
  & Closest          & \underline{93.29} & \underline{93.66} & \underline{88.52} & \underline{1.72} & 48.25 \\
  & Vote             & 85.10 & \underline{93.66} & 77.17 & \underline{0.84} & 48.19 \\
\bottomrule
\end{tabular}
\end{table}

\subsection{WavLM X-Vector (512-dim)}
\label{sec:results_wavlm_xvec}

\begin{table}[H]
\centering
\caption{WavLM X-Vector results (cosine, $C{=}1$, schema: \code{wavlm\_xvector}).}
\label{tab:results_wavlm_xvec}
\setlength{\tabcolsep}{4pt}
\begin{tabular}{llrrrrr}
\toprule
\textbf{Architecture} & \textbf{Aggregation} &
\textbf{Top-1\%} & \textbf{Hit@10\%} &
\textbf{F1\%} & \textbf{EER\%} & \textbf{ms} \\
\midrule
\multirow{3}{*}{Brute-force}
  & Speaker centroid & 95.84 & 99.10 & 92.55 & 14.97 & 0.051 \\
  & Closest          & \textbf{97.42} & \textbf{99.65} & \textbf{94.96} & 17.86 & 238.34 \\
  & Vote             & 96.83 & \textbf{99.65} & 94.03 & \textbf{14.18} & 237.26 \\
\midrule
\multirow{3}{*}{Metadata hier.}
  & Speaker centroid & 32.83 & 33.62 & 28.50 & 4.18 & 0.038 \\
  & Closest          & 33.07 & 33.63 & 28.59 & 3.70 & 26.24 \\
  & Vote             & 33.10 & \underline{33.63} & 28.74 & \underline{3.52} & 26.47 \\
\midrule
\multirow{3}{*}{$K$-Means}
  & Speaker centroid & \underline{93.44} & \underline{95.78} & \underline{88.94} & 11.78 & \textbf{0.023} \\
  & Closest          & \underline{93.04} & \underline{95.95} & \underline{88.18} & 13.45 & 14.01 \\
  & Vote             & 81.13 & \underline{95.95} & 73.12 & \underline{7.83} & 14.04 \\
\bottomrule
\end{tabular}
\end{table}

\subsection{WavLM Mean-Pooled (768-dim)}
\label{sec:results_wavlm}

\begin{table}[H]
\centering
\caption{WavLM mean-pooled results (cosine, $C{=}1$, schema: \code{wavlm}).}
\label{tab:results_wavlm}
\setlength{\tabcolsep}{4pt}
\begin{tabular}{llrrrrr}
\toprule
\textbf{Architecture} & \textbf{Aggregation} &
\textbf{Top-1\%} & \textbf{Hit@10\%} &
\textbf{F1\%} & \textbf{EER\%} & \textbf{ms} \\
\midrule
\multirow{3}{*}{Brute-force}
  & Speaker centroid & 26.19 & 49.87 & 18.01 & 38.78 & 0.039 \\
  & Closest          & 26.96 & 51.90 & 19.01 & 33.86 & 259.54 \\
  & Vote             & \textbf{29.66} & \textbf{53.08} & \textbf{20.44} & 35.64 & 258.10 \\
\midrule
\multirow{3}{*}{Metadata hier.}
  & Speaker centroid & 7.46 & 12.79 & 4.99 & 27.20 & 0.032 \\
  & Closest          & 7.04 & 12.49 & 5.20 & \textbf{25.21} & 21.14 \\
  & Vote             & 7.35 & 12.66 & 4.89 & 26.49 & 21.13 \\
\midrule
\multirow{3}{*}{$K$-Means}
  & Speaker centroid & \underline{23.62} & \underline{43.11} & \underline{16.39} & 37.98 & \textbf{0.081} \\
  & Closest          & 17.95 & 34.64 & 12.72 & 32.63 & 14.45 \\
  & Vote             & 14.20 & 35.18 & 9.02 & 31.17 & 14.43 \\
\bottomrule
\end{tabular}
\end{table}

\subsection{Metadata Hierarchy: Varying Probe Depth}
\label{sec:results_probes}

The following tables extend the metadata hierarchy results by varying
the nationality probe depth $P \in \{2, 5, 10\}$.
$P{=}1$ results appear in the per-model tables above.

\begin{table}[H]
\centering
\caption{ECAPA-TDNN --- metadata hierarchy with varying probe depth $P$.}
\label{tab:probes_ecapa}
\setlength{\tabcolsep}{4pt}
\begin{tabular}{rlrrrrr}
\toprule
$P$ & \textbf{Aggregation} &
\textbf{Top-1\%} & \textbf{Hit@10\%} &
\textbf{F1\%} & \textbf{EER\%} & \textbf{ms} \\
\midrule
\multirow{3}{*}{2}
  & Speaker centroid & 47.71 & 47.78 & 41.31 & \textbf{0.15} & \textbf{0.044} \\
  & Closest          & 47.75 & 47.79 & 41.84 & 0.39 & 32.79 \\
  & Vote             & 47.73 & 47.79 & 41.85 & 0.44 & 33.00 \\
\midrule
\multirow{3}{*}{5}
  & Speaker centroid & 65.86 & 65.99 & 56.45 & \textbf{0.21} & \textbf{0.072} \\
  & Closest          & 65.91 & 66.00 & 57.39 & 0.40 & 33.36 \\
  & Vote             & 65.89 & 66.00 & 57.43 & 0.40 & 33.52 \\
\midrule
\multirow{3}{*}{10}
  & Speaker centroid & 73.87 & 74.08 & 64.82 & \textbf{0.26} & \textbf{0.097} \\
  & Closest          & \textbf{73.94} & \textbf{74.09} & \textbf{65.91} & 0.44 & 30.49 \\
  & Vote             & 73.93 & \textbf{74.09} & 65.86 & 0.44 & 30.54 \\
\bottomrule
\end{tabular}
\end{table}

\begin{table}[H]
\centering
\caption{WavLM X-Vector --- metadata hierarchy with varying probe depth $P$.}
\label{tab:probes_wavlm_xvec}
\setlength{\tabcolsep}{4pt}
\begin{tabular}{rlrrrrr}
\toprule
$P$ & \textbf{Aggregation} &
\textbf{Top-1\%} & \textbf{Hit@10\%} &
\textbf{F1\%} & \textbf{EER\%} & \textbf{ms} \\
\midrule
\multirow{3}{*}{2}
  & Speaker centroid & 57.38 & 58.96 & 50.44 & 4.62 & \textbf{0.058} \\
  & Closest          & 57.92 & 58.97 & 50.85 & 4.28 & 45.75 \\
  & Vote             & 57.77 & 58.97 & 50.76 & \textbf{3.83} & 45.55 \\
\midrule
\multirow{3}{*}{5}
  & Speaker centroid & 80.97 & 83.36 & 74.10 & 7.69 & \textbf{0.082} \\
  & Closest          & 81.79 & 83.49 & 75.29 & 7.27 & 67.58 \\
  & Vote             & 81.55 & 83.49 & 75.06 & \textbf{6.12} & 67.45 \\
\midrule
\multirow{3}{*}{10}
  & Speaker centroid & 89.53 & 92.30 & 83.44 & 10.46 & \textbf{0.110} \\
  & Closest          & \textbf{90.56} & \textbf{92.47} & \textbf{84.77} & 10.26 & 82.46 \\
  & Vote             & 90.24 & \textbf{92.47} & 84.47 & \textbf{8.85} & 82.34 \\
\bottomrule
\end{tabular}
\end{table}

\begin{table}[H]
\centering
\caption{WavLM Mean-Pooled --- metadata hierarchy with varying probe depth $P$.}
\label{tab:probes_wavlm}
\setlength{\tabcolsep}{4pt}
\begin{tabular}{rlrrrrr}
\toprule
$P$ & \textbf{Aggregation} &
\textbf{Top-1\%} & \textbf{Hit@10\%} &
\textbf{F1\%} & \textbf{EER\%} & \textbf{ms} \\
\midrule
\multirow{3}{*}{2}
  & Speaker centroid & 12.29 & 22.71 & 8.88  & 30.67 & \textbf{0.058} \\
  & Closest          & 12.13 & 22.71 & 9.27  & \textbf{27.85} & 35.80 \\
  & Vote             & 12.75 & 23.11 & 9.09  & 27.96 & 35.69 \\
\midrule
\multirow{3}{*}{5}
  & Speaker centroid & 21.35 & 40.69 & 15.04 & 36.02 & \textbf{0.159} \\
  & Closest          & 21.09 & 41.04 & 15.05 & \textbf{31.58} & 73.75 \\
  & Vote             & \textbf{22.76} & \textbf{41.94} & \textbf{15.86} & 33.01 & 74.05 \\
\midrule
\multirow{3}{*}{10}
  & Speaker centroid & 24.05 & 46.62 & 16.61 & 37.86 & \textbf{0.285} \\
  & Closest          & 24.29 & 47.36 & 17.17 & 33.08 & 114.96 \\
  & Vote             & \textbf{26.34} & \textbf{48.37} & \textbf{17.98} & 34.58 & 116.52 \\
\bottomrule
\end{tabular}
\end{table}

\subsection{Cross-Model Comparison}
\label{sec:results_crossmodel}

Table~\ref{tab:crossmodel} places representative operating points of
all three models side by side.

\begin{table}[H]
\centering
\caption{Cross-model comparison at selected operating points.}
\label{tab:crossmodel}
\setlength{\tabcolsep}{4pt}
\begin{tabular}{llrrr}
\toprule
\textbf{Configuration} & \textbf{Metric} &
\textbf{ECAPA} & \textbf{WavLM-XV} & \textbf{WavLM} \\
\midrule
\multirow{3}{*}{BF + closest}
  & Top-1\%  & \textbf{99.75} & 97.42  & 26.96  \\
  & EER\%    & 0.79  & 17.86  & 33.86  \\
  & ms/query & 37.81 & 238.34 & 259.54 \\
\midrule
\multirow{3}{*}{BF + centroid}
  & Top-1\%  & \textbf{99.64} & 95.84  & 26.19  \\
  & EER\%    & \textbf{0.61}  & 14.97  & 38.78  \\
  & ms/query & 0.030 & 0.051  & 0.039  \\
\midrule
\multirow{3}{*}{K-Means + centroid}
  & Top-1\%  & \textbf{93.56} & 93.44  & 23.62  \\
  & EER\%    & 0.90  & 11.78  & 37.98  \\
  & ms/query & \textbf{0.015} & 0.023  & 0.081  \\
\midrule
\multirow{3}{*}{K-Means + closest}
  & Top-1\%  & \textbf{93.29} & 93.04  & 17.95  \\
  & EER\%    & 1.72  & 13.45  & 32.63  \\
  & ms/query & 48.25 & 14.01  & 14.45  \\
\bottomrule
\end{tabular}
\end{table}

% ============================================================
\section{Analysis}
\label{sec:analysis}

\subsection{Effect of Search Architecture}
\label{sec:analysis_arch}

\paragraph{Brute-force.}
For ECAPA-TDNN, brute-force delivers near-perfect Top-1 accuracy
(99.75\% closest, 99.64\% centroid) with EERs below 1\%.
The minor gap between \emph{closest} and \emph{centroid} ($<$0.12\,pp)
suggests that averaging utterances per speaker is a near-lossless
compression of identity information when embeddings are well-trained.
For WavLM X-Vector the gap is larger (97.42\% vs.\ 95.84\%),
indicating slightly less intra-speaker consistency in the 512-dim space.
For WavLM mean-pooled, brute-force itself cannot exceed 29.66\%,
confirming that embedding quality is the binding constraint: no search
strategy can recover from poorly discriminative representations.

\paragraph{Metadata hierarchy.}
Results collapse to ${\approx}33\%$ for ECAPA and WavLM X-Vector
and to ${\approx}7\%$ for WavLM mean-pooled.
The random-chance baseline for 1,251 speakers is
$1/1251 \approx 0.08\%$, so the 33\% figure is non-trivial
yet far below any practical operating threshold.
The root cause is coarse geographic partitioning: a single nationality
may cover hundreds of speakers while another covers only a handful,
creating highly imbalanced candidate pools.
A routing error in the nationality selection step is unrecoverable
because the filtered candidate pool then excludes the true speaker.
Hard metadata gating is therefore contraindicated unless partition
balance and routing accuracy are independently verified.

\paragraph{Effect of metadata probe depth.}
Tables~\ref{tab:probes_ecapa}--\ref{tab:probes_wavlm} show that
increasing the nationality probe depth $P$ steadily recovers accuracy
for all models, but the rate and ceiling differ sharply.
For ECAPA-TDNN, Top-1 rises from 33.3\% at $P{=}1$ to 47.8\%,
65.9\%, and 73.9\% at $P{=}2$, 5, and 10 respectively, while EER
remains below 0.45\% throughout, confirming that the embeddings
themselves are well-calibrated and the loss at $P{=}1$ is purely a
routing coverage problem.
For WavLM X-Vector the gains are larger: Top-1 jumps from 33.1\%
to 57.9\% ($P{=}2$), 81.8\% ($P{=}5$), and 90.6\% ($P{=}10$),
approaching brute-force closest (97.4\%) at $P{=}10$ with the closest
aggregation.
This shows that WavLM X-Vector embeddings are discriminative enough
for metadata-guided search given a sufficiently wide candidate pool.
For WavLM mean-pooled, even at $P{=}10$ Top-1 reaches only 26.3\%
(vote), barely above the brute-force baseline of 29.7\%, confirming
that the bottleneck is the embedding quality rather than the routing
breadth.
In all cases, the speaker-centroid aggregation preserves near-zero
latency ($<$0.3\,ms) regardless of $P$, while closest and vote
aggregation scale linearly with the enlarged candidate pool.

\paragraph{K-means cluster routing.}
For ECAPA, K-means centroid reaches 93.56\% Top-1 at 0.015\,ms/query
, a $2{,}520{\times}$ speed-up over brute-force closest (37.81\,ms)
with a 6.19\,pp accuracy cost.
The accuracy loss is attributable to routing errors: speakers who
straddle cluster boundaries have their utterances split across clusters,
and a single-probe search ($C{=}1$) misses any utterance not in the
queried cluster.
For WavLM X-Vector, K-means centroid achieves 93.44\% at 0.023\,ms,
nearly matching ECAPA despite the larger embedding dimension.
For WavLM mean-pooled, K-means degrades below brute-force
(23.62\% vs.\ 26.19\%), indicating that the cluster structure in
this space does not correlate with speaker identity.

\subsection{Effect of Aggregation Mode}
\label{sec:analysis_agg}

\paragraph{Speaker centroid vs.\ closest.}
Under brute-force, \emph{closest} and \emph{centroid} differ
by at most 0.12\,pp (ECAPA) and 1.58\,pp (WavLM X-Vector) in Top-1.
Yet centroid queries require no DB round-trip for ranking: the cosine
dot-product of the query against 1,251 pre-computed per-speaker mean
vectors runs entirely in Python/NumPy, making centroid queries roughly
$1{,}260{\times}$ faster than closest in the brute-force case
(0.030\,ms vs.\ 37.81\,ms).
This makes \emph{brute-force centroid} the dominant strategy for
offline identification: near-optimal accuracy at negligible latency.

\paragraph{Vote aggregation.}
Aggregating the top-25 candidate utterances by majority vote degrades
accuracy for K-means in all three models.
With $C{=}1$ and an average cluster of ${\approx}600$ utterances,
the candidate pool is dominated by speakers with many utterances in
the winning cluster.
Vote aggregation amplifies this positional bias, producing results
consistently worse than simply selecting the closest utterance.
Under brute-force, vote also underperforms closest for ECAPA
(99.68\% vs.\ 99.75\%) and substantially elevates EER (3.57\% vs.\
0.79\%), suggesting that distant false-positive candidates in the
top-25 corrupt the vote signal.

\subsection{Accuracy--Latency Trade-off}
\label{sec:analysis_tradeoff}

\begin{table}[H]
\centering
\caption{Key operating points across all models and architectures.
  Speedup relative to ECAPA brute-force closest (37.81\,ms).}
\label{tab:tradeoff}
\setlength{\tabcolsep}{4pt}
\begin{tabular}{llrrrr}
\toprule
\textbf{Model} & \textbf{Configuration} &
\textbf{Top-1\%} & \textbf{EER\%} &
\textbf{ms} & \textbf{Speedup} \\
\midrule
ECAPA    & BF + closest           & 99.75 & 0.79  & 37.81  & $1{\times}$       \\
ECAPA    & BF + centroid          & 99.64 & 0.61  & 0.030  & $1{,}260{\times}$ \\
ECAPA    & K-Means + centroid     & 93.56 & 0.90  & 0.015  & $2{,}520{\times}$ \\
ECAPA    & K-Means + closest      & 93.29 & 1.72  & 48.25  & $0.8{\times}$     \\
\midrule
WavLM-XV & BF + closest           & 97.42 & 17.86 & 238.34 & $0.16{\times}$    \\
WavLM-XV & BF + centroid          & 95.84 & 14.97 & 0.051  & $742{\times}$     \\
WavLM-XV & K-Means + centroid     & 93.44 & 11.78 & 0.023  & $1{,}644{\times}$ \\
\midrule
WavLM    & BF + vote (best)       & 29.66 & 35.64 & 258.10 &,               \\
WavLM    & BF + centroid          & 26.19 & 38.78 & 0.039  &,               \\
\bottomrule
\end{tabular}
\end{table}

The ECAPA K-Means centroid configuration (0.015\,ms, 93.56\%)
represents the practical optimum for a scalable real-time system:
a $2{,}520{\times}$ latency reduction at a cost of only 6.19\,pp
Top-1 and 0.29\,pp EER relative to brute-force closest.
The brute-force centroid is the better choice for offline batch
processing (0.030\,ms, 99.64\%), where a DB query per utterance is
not required and the accuracy loss is negligible.

WavLM X-Vector K-Means centroid achieves nearly identical Top-1
(93.44\%) but at a substantially higher EER (11.78\% vs.\ 0.90\%)
and slightly longer latency (0.023\,ms vs.\ 0.015\,ms), making
ECAPA the superior model on every dimension.

\subsection{Embedding Model Comparison}
\label{sec:analysis_model}

The results reveal a clear three-tier hierarchy:

\begin{enumerate}
  \item \textbf{ECAPA-TDNN}, best on all accuracy and EER metrics,
        lowest latency for centroid search due to the smallest
        embedding dimension (192).
        Its task-specific AAM-softmax fine-tuning produces
        well-calibrated cosine similarity scores, reflected in EER
        below 1\% even for K-means routing.
  \item \textbf{WavLM X-Vector}, 97.42\% Top-1 under brute-force
        closest, but suffers from high EER ($\geq$14\%) across all
        configurations.
        This suggests that the similarity score distribution is
        poorly calibrated: the model ranks the correct speaker first
        in most cases but assigns too-similar scores to impostors,
        making it unreliable for threshold-based accept/reject decisions.
        The higher latency for closest search (238\,ms vs.\ 37.81\,ms
        for ECAPA) reflects the larger 512-dim vectors.
  \item \textbf{WavLM mean-pooled}, completely unsuitable for
        closed-set speaker identification without fine-tuning.
        Best Top-1 is 29.66\% (vote, brute-force), EER exceeds 33\%
        in all configurations, and K-means routing actively hurts
        performance relative to brute-force.
        Generic self-supervised features encode phonetic, prosodic,
        and linguistic content that is not speaker-discriminative at
        the identity level required for this task.
\end{enumerate}

The key takeaway is that \emph{the training objective of the embedding
model dominates all other architectural choices}.
Switching from ECAPA to WavLM mean-pooled costs $>70$\,pp Top-1,
while switching from brute-force to K-means routing within the same
model costs $<7$\,pp.
Practitioners should therefore prioritise embedding model quality
before optimising search strategy.

% ============================================================
\section{Conclusions}
\label{sec:conclusion}

We evaluated 27 experimental configurations (3 models $\times$
3 search architectures $\times$ 3 aggregation modes) on VoxCeleb1
with 8,676 test utterances against 1,251 enrolled speakers.
The principal findings are:

\begin{enumerate}
  \item \textbf{Embedding model training objective is paramount.}
        Fine-tuned speaker verification embeddings (ECAPA-TDNN)
        outperform generic self-supervised representations (WavLM
        mean-pooled) by $>70$\,pp Top-1 regardless of search strategy.
  \item \textbf{Brute-force centroid is the practical optimum for ECAPA.}
        At 0.030\,ms/query and 99.64\% Top-1, it eliminates DB
        queries for ranking entirely and loses only 0.11\,pp versus
        the accuracy ceiling.
  \item \textbf{K-Means centroid is the scalable real-time option.}
        0.015\,ms, 93.56\% Top-1, EER 0.90\%, a $2{,}520{\times}$
        speed-up suitable for latency-sensitive applications.
  \item \textbf{Metadata hierarchy is not viable as a hard filter.}
        ${\approx}33\%$ Top-1 for all models reflects imbalanced
        nationality partitions rather than genuine speaker similarity.
  \item \textbf{Vote aggregation degrades under K-Means routing.}
        The limited cluster candidate pool biases votes towards
        majority-utterance speakers, reducing accuracy by up to
        8.5\,pp relative to K-Means closest.
  \item \textbf{WavLM X-Vector EER is problematic.}
        Despite 97.42\% Top-1, an EER of 17.9\% makes this model
        unreliable for any application requiring threshold-based
        accept/reject decisions.
\end{enumerate}

Future work should explore multi-probe K-Means ($C{>}1$) to close the
accuracy gap with brute-force, HNSW approximate search as an
intermediate operating point, and AAM-softmax fine-tuning of WavLM
to evaluate whether the EER penalty can be eliminated while retaining
the richer self-supervised representations.

% ============================================================
\clearpage
\appendix
\section{Database Schema SQL}
\label{app:schema}

\begin{lstlisting}[language=SQL,
  caption={Schema template (data\_base/schema.template.sql).}]
CREATE SCHEMA IF NOT EXISTS {schema},
SET search_path TO {schema},

CREATE TABLE speakers (
    speaker_id  VARCHAR(20) PRIMARY KEY,
    gender      VARCHAR(10) NOT NULL,
    nationality VARCHAR(20) NOT NULL
),
CREATE INDEX idx_spk_gn ON speakers(gender, nationality),

CREATE TABLE audio_embeddings (
    id         INT GENERATED ALWAYS AS IDENTITY PRIMARY KEY,
    speaker_id VARCHAR(20) REFERENCES speakers(speaker_id),
    file_path  TEXT NOT NULL,
    cluster_id INT,
    embedding  VECTOR({dim}) NOT NULL
),
CREATE INDEX idx_ae_spk  ON audio_embeddings(speaker_id),
CREATE INDEX idx_ae_cl   ON audio_embeddings(cluster_id),
CREATE INDEX idx_ae_hnsw ON audio_embeddings
    USING hnsw(embedding vector_cosine_ops)
    WITH (m=16, ef_construction=64),

CREATE TABLE metadata_centroids (
    category        VARCHAR(20) NOT NULL,
    category_value  VARCHAR(50) NOT NULL,
    centroid_vector VECTOR({dim}) NOT NULL,
    PRIMARY KEY (category, category_value)
),

CREATE TABLE cluster_centroids (
    cluster_id      INT PRIMARY KEY,
    centroid_vector VECTOR({dim}) NOT NULL
),
ALTER TABLE audio_embeddings ADD CONSTRAINT fk_cl
    FOREIGN KEY (cluster_id)
    REFERENCES cluster_centroids(cluster_id),
\end{lstlisting}

% ============================================================
\clearpage
\begin{thebibliography}{99}

\bibitem{snyder2018xvector}
D.\ Snyder, D.\ Garcia-Romero, G.\ Sell, D.\ Povey, and S.\ Khudanpur,
``$X$-vectors: Robust DNN embeddings for speaker recognition,''
in \textit{Proc.\ IEEE ICASSP}, 2018, pp.~5329--5333.

\bibitem{desplanques2020ecapa}
B.\ Desplanques, J.\ Thienpondt, and K.\ Demuynck,
``ECAPA-TDNN: Emphasized channel attention, propagation and aggregation
in TDNN-based speaker verification,''
in \textit{Proc.\ INTERSPEECH}, 2020, pp.~3830--3834.

\bibitem{chen2022wavlm}
S.\ Chen \textit{et al.},
``WavLM: Large-scale self-supervised pre-training for full stack speech processing,''
\textit{IEEE J.\ Sel.\ Topics Signal Process.}, vol.~16, no.~6,
pp.~1505--1518, 2022.

\bibitem{nagrani2017voxceleb}
A.\ Nagrani, J.\ S.\ Chung, and A.\ Zisserman,
``VoxCeleb: A large-scale speaker identification dataset,''
in \textit{Proc.\ INTERSPEECH}, 2017, pp.~2616--2620.

\bibitem{ravanelli2021speechbrain}
M.\ Ravanelli \textit{et al.},
``SpeechBrain: A general-purpose speech toolkit,''
\textit{arXiv:2106.04624}, 2021.


\bibitem{pgvector2023}
A.\ Kane,
``pgvector: Open-source vector similarity search for PostgreSQL,''
\url{https://github.com/pgvector/pgvector}, 2023.

\end{thebibliography}

\end{document}